% !PRINTED 41
\weberpage{41}

We distinguish four cases:

1. $f''''(\alpha)$ and $f(\alpha)$ not $=0$ [i.\,e. $f(x)$ of the fourth degree and $\alpha$ not a root of $f(x)$].

In this case formula (2) for $n=0$ allows $S_3$ to be expressed by $S_0, S_1, S_2$, and hence $S_n$ for every positive $n$ by $S_0, S_1, S_2$. If, however, one puts $n=-1,-2,\ldots$, one obtains $S_{-2}, S_{-3}, \ldots$ expressed by $S_{-1}, S_0, S_1, S_2$. In this case, therefore, the four fundamental integrals remain
\[
S_{-1}, \quad S_0, \quad S_1, \quad S_2.
\]

2. $f''''(\alpha)$ not $=0$, $f(\alpha)=0$. In this case $S_{-1}$ can still be expressed by $S_0, S_1, S_2$, and one obtains as fundamental integrals
\[
S_0, \quad S_1, \quad S_2.
\]

3. $f''''(\alpha)=0$, $f(\alpha)$ not $=0$. Here the fundamental integrals are
\[
S_{-1}, \quad S_0, \quad S_1.
\]

4. $f''''(\alpha)=0$, $f(\alpha)=0$. Here the fundamental integrals are
\[
S_0, \quad S_1.
\]

The result of this is therefore:

I. If $f(x)$ is of the fourth degree ($a_0=1$), then all integrals $\Omega(x)$ can be expressed by integrals of the form
\[
S_0 = \int \frac{dx}{\sqrt{f(x)}}, \quad S_1 = \int \frac{x\,dx}{\sqrt{f(x)}}, \quad S_2 = \int \frac{x^2\,dx}{\sqrt{f(x)}},
\]
\[
S_{-1} = \int \frac{dx}{(x-\alpha)\sqrt{f(x)}},
\]
in which $S_{-1}$ still depends on the parameter $\alpha$.

II. If, however, $f(x)$ is of the third degree, then the three integrals suffice
\[
S_0 = \int \frac{dx}{\sqrt{f(x)}}, \quad S_1 = \int \frac{x\,dx}{\sqrt{f(x)}}, \quad S_{-1} = \int \frac{dx}{(x-\alpha)\sqrt{f(x)}}.
\]

In both cases $S_0$ is \gesperrt{of the first kind}. If $f(x)$ is of the third degree, then $S_1$ is of the second and $S_{-1}$ of the third kind. If on the other hand $f(x)$ is of the fourth degree, then $S_1$, $S_2$ and $S_{-1}$ are of the third kind. But one can compose from $S_1$ and $S_2$ an integral of the second kind. For we have
%CONT>>
